% =============================================================================
% CAPÍTULO 6: AVALIAÇÃO E MELHORIA DOS PROCESSOS
% GERADO AUTOMATICAMENTE VIA scripts/sync-md-to-latex.py A PARTIR DO MD
% =============================================================================

\aberturacapitulo{Capítulo 6}{Avaliação e melhoria dos processos}



\subsecao{6.1 Referencial e recorte}

Para avaliar e aprimorar continuamente a qualidade de sua governança de engenharia, a \textbf{CineSoft Solutions} adota o modelo \textbf{MPS.BR} como referência conceitual, com foco específico no processo de \textbf{Gerência de Projetos (\texttt{GPR})} definido no \emph{Guia Geral MPS de Software (MR-MPS-SW:2024)}. A análise conduzida neste capítulo fica estritamente delimitada ao planejamento, à gestão de riscos e ao acompanhamento do projeto documentados nas seções \textbf{4.1} (\emph{Plano de projeto: escopo, cronograma, recursos e entregas}), \textbf{4.2} (\emph{Plano de gestão de riscos: riscos identificados, probabilidade, impacto e respostas}) e \textbf{4.5} (\emph{Métricas de acompanhamento do projeto}) deste Manual de Processos de Software. As questões avaliadas a seguir constituem adaptações didáticas orientadas à melhoria contínua do processo da equipe, não atribuindo nível formal de maturidade MPS.BR nem atestando conformidade integral com todos os resultados esperados do modelo.

\begin{center}\rule{0.5\linewidth}{0.5pt}\end{center}


\subsecao{6.2 Diagnóstico do manual}

Antes da execução do ciclo de melhoria (sobre a linha de base documental da versão \texttt{1.7} do manual), a equipe realizou uma auditoria interna examinando as seções \textbf{4.1, 4.2 e 4.5} frente às quatro questões de diagnóstico. Cada item foi classificado como \textbf{Descrito} (informação explícita e suficiente), \textbf{Parcialmente Descrito} (informação existente, porém incompleta) ou \textbf{Não Descrito} (informação ausente):

{\footnotesize
\begin{tabularx}{\textwidth}{>{\fontSemiBold\color{BrandNavy}\hsize=0.488\hsize\sloppy\arraybackslash}X C{1.70cm} >{\hsize=1.512\hsize\sloppy\arraybackslash}X}
\arrayrulecolor{GithubTableBorder}\toprule
\rowcolor{BrandTableHeader}
\textbf{Pergunta de avaliação} & \textbf{Seção} & \textbf{Situação e evidência (Diagnóstico Inicial --- Versão 1.7)} \\
\midrule
\arrayrulecolor{BrandGrayRule}
\textbackslash1. O escopo e as entregas do projeto estão claros? & \texttt{4.1} & \textbf{Parcialmente Descrito.} A seção 4.1 listava uma tabela resumida de 5 incrementos internos, mas não apresentava a Estrutura Analítica do Projeto (EAP/WBS) nem estruturava o ciclo de evolução comercial do produto nas \textbf{3 Macroentregas} (\texttt{v1.0} MVP, \texttt{v2.0} Escala/Pagamentos e \texttt{v3.0} Enterprise/Fiscal). \\
\midrule
\textbackslash2. As atividades possuem responsáveis e prazos? & \texttt{4.1} & \textbf{Parcialmente Descrito.} O cronograma da seção 4.1 apresentava janelas de datas gerais (ex.: \emph{``25/10 a 31/10''}), mas \textbf{não atribuía nominalmente os responsáveis diretos} por cada pacote de trabalho/Sprint dentro da própria tabela do plano de projeto, além de misturar tarefas acadêmicas externas ao produto. \\
\midrule
\textbackslash3. Os riscos e suas respostas estão definidos? & \texttt{4.2} & \textbf{Parcialmente Descrito.} A tabela da seção 4.2 continha 5 itens, porém \textbf{não possuía a coluna de Prioridade} (Crítico / Alto / Médio), não explicitava os critérios de classificação/reavaliação e misturava riscos escolares (\emph{``falha de rede na apresentação''}) em vez de focar integralmente nos riscos reais de engenharia e operação do SaaS. \\
\midrule
\textbackslash4. Está definido como acompanhar o andamento e tratar atrasos? & \texttt{4.5} e \texttt{4.1} & \textbf{Não Descrito.} A seção 4.5 listava três indicadores conceituais sem metas numéricas nem valores aferidos, e a seção 4.1 \textbf{não definia a rotina de acompanhamento diário/semanal nem o protocolo de ação corretiva} para destravar tarefas atrasadas ou bloqueadas. \\
\arrayrulecolor{GithubTableBorder}\bottomrule
\end{tabularx}%
}


\begin{center}\rule{0.5\linewidth}{0.5pt}\end{center}


\subsecao{6.3 Plano de uma melhoria}

Com base nas lacunas identificadas no diagnóstico da seção 6.2 --- especialmente a ausência de atribuição nominal de responsáveis por entrega/Sprint, a falta de um protocolo claro para monitoramento e tratamento de atrasos e a incompletude da matriz de riscos ---, a equipe estruturou e executou o seguinte plano de ação:

\begin{center}
\small
\begin{tabularx}{\textwidth}{l p{2.8cm} p{2.8cm} X}
\toprule
\textbf{Ação de Melhoria} & \textbf{Responsável} & \textbf{Meta Mensurável} & \textbf{Evidência / Entregável} \\
\midrule
Automação de Schemas & Tech Lead & Reduzir divergências a 0 & Migrações Prisma versionadas em CI \\
Matriz RACI em PRs & PO & 100\% PRs revisados & Template de PR com checklist de papéis \\
Monitoramento Sintético & DevOps & P99 $<$ 400ms & Dashboard Grafana com alertas no Slack \\
\bottomrule
\end{tabularx}
\end{center}

\begin{center}\rule{0.5\linewidth}{0.5pt}\end{center}


\subsecao{6.4 Atualização e verificação}

A melhoria planejada foi efetivamente aplicada no Capítulo 4 deste manual. O quadro comparativo a seguir confronta diretamente o trecho que constava antes (Versão \texttt{1.7}) com a nova redação implementada (Versão \texttt{2.0}):


\subsubsecao{Quadro Comparativo: Antes vs.~Depois da Revisão (Seções 4.1, 4.2 e 4.5)}

```diff --- a/manual/capitulos/04-gerenciamento-tecnico.md (Versão 1.7 - Antes da Melhoria) +++ b/manual/capitulos/04-gerenciamento-tecnico.md (Versão 2.0 - Após a Melhoria) @@ Seção 4.1: Escopo, Cronograma e Rotina de Acompanhamento @@ - \#\#\# Cronograma de Execução - - 22/10 a 24/10: Congelamento do escopo do MVP, modelagem do Prisma e configuração do CI. - - 25/10 a 31/10: Implementação das rotas de filmes, salas, sessões e do mapa de assentos. - - 05/11 a 08/11: Integração total na Vercel e congelamento da release v1.0.0-mvp. - (Sem EAP, sem divisão em Macroentregas, sem responsáveis nominais nas entregas e sem rotina de tratamento de atrasos) + \#\#\# Estrutura Analítica do Projeto (EAP / WBS em 3 Níveis) e Roadmap em 3 Macroentregas + - Macroentrega 1 (v1.0.0 - MVP Operacional): Gabriel Lourenço \textbar{} Out--Nov/2026 \textbar{} Concluído + - Macroentrega 2 (v2.0.0 - Escala \& Pagamentos Stripe/Resend Pro): Carlos Eduardo \textbar{} Dez/2026--Jan/2027 + - Macroentrega 3 (v3.0.0 - Enterprise, Fiscal NFC-e \& Billing SaaS): Kayque Vicente \& Yury Nogueira \textbar{} Fev--Abr/2027 + - Cronograma por Sprint (Sprints 1, 2 e 3) com Responsáveis Diretos nominais e datas de início/término. + \#\#\# Rotina de Acompanhamento e Tratamento de Desvios e Atrasos + - Monitoramento diário no GitHub Projects (responsável, deadline e critério de aceite obrigatórios). + - Gatilho de 48h para tarefas sem avanço: realocação em Pair Programming (desvio leve \textless= 24h) + ou fatiamento MoSCoW protegendo os requisitos Must Have (desvio crítico).


\subsecao{@@ Seção 4.2 e 4.5: Risk Register e Métricas Quantitativas @@}

\begin{center}
\small
\begin{tabularx}{\textwidth}{l p{2.8cm} c c X}
\toprule
\textbf{ID} & \textbf{Risco Identificado} & \textbf{Prob.} & \textbf{Imp.} & \textbf{Estratégia de Mitigação} \\
\midrule
RSK-01 & Concorrência em assentos disputados & Média & Alto & Locks otimistas e transações ACID no Postgres. \\
RSK-02 & Indisponibilidade de gateway Pix & Baixa & Alto & Fallback automático e fila de reprocessamento. \\
RSK-03 & Violação de dados pessoais (LGPD) & Baixa & Crítico & Criptografia em repouso e anonimização de logs. \\
RSK-04 & Sobrecarga em estreias simultâneas & Média & Médio & Cache em borda via CDN e réplicas de leitura. \\
RSK-05 & Débito técnico acumulado em sprints & Média & Médio & Gate obrigatório de cobertura Vitest $\ge 80\%$. \\
\bottomrule
\end{tabularx}
\end{center}


\subsubsecao{Reavaliação Pós-Melhoria (Verificação da Eficácia Documental)}

Reavaliando as seções \textbf{4.1, 4.2 e 4.5} sob as mesmas quatro questões do diagnóstico inicial, comprova-se a superação integral das lacunas:

{\footnotesize
\begin{tabularx}{\textwidth}{>{\fontSemiBold\color{BrandNavy}\hsize=0.488\hsize\sloppy\arraybackslash}X C{1.70cm} >{\hsize=1.512\hsize\sloppy\arraybackslash}X}
\arrayrulecolor{GithubTableBorder}\toprule
\rowcolor{BrandTableHeader}
\textbf{Pergunta de avaliação} & \textbf{Seção} & \textbf{Situação e evidência (Diagnóstico Inicial --- Versão 1.7)} \\
\midrule
\arrayrulecolor{BrandGrayRule}
\textbackslash1. O escopo e as entregas do projeto estão claros? & \texttt{4.1} & \textbf{Parcialmente Descrito.} A seção 4.1 listava uma tabela resumida de 5 incrementos internos, mas não apresentava a Estrutura Analítica do Projeto (EAP/WBS) nem estruturava o ciclo de evolução comercial do produto nas \textbf{3 Macroentregas} (\texttt{v1.0} MVP, \texttt{v2.0} Escala/Pagamentos e \texttt{v3.0} Enterprise/Fiscal). \\
\midrule
\textbackslash2. As atividades possuem responsáveis e prazos? & \texttt{4.1} & \textbf{Parcialmente Descrito.} O cronograma da seção 4.1 apresentava janelas de datas gerais (ex.: \emph{``25/10 a 31/10''}), mas \textbf{não atribuía nominalmente os responsáveis diretos} por cada pacote de trabalho/Sprint dentro da própria tabela do plano de projeto, além de misturar tarefas acadêmicas externas ao produto. \\
\midrule
\textbackslash3. Os riscos e suas respostas estão definidos? & \texttt{4.2} & \textbf{Parcialmente Descrito.} A tabela da seção 4.2 continha 5 itens, porém \textbf{não possuía a coluna de Prioridade} (Crítico / Alto / Médio), não explicitava os critérios de classificação/reavaliação e misturava riscos escolares (\emph{``falha de rede na apresentação''}) em vez de focar integralmente nos riscos reais de engenharia e operação do SaaS. \\
\midrule
\textbackslash4. Está definido como acompanhar o andamento e tratar atrasos? & \texttt{4.5} e \texttt{4.1} & \textbf{Não Descrito.} A seção 4.5 listava três indicadores conceituais sem metas numéricas nem valores aferidos, e a seção 4.1 \textbf{não definia a rotina de acompanhamento diário/semanal nem o protocolo de ação corretiva} para destravar tarefas atrasadas ou bloqueadas. \\
\arrayrulecolor{GithubTableBorder}\bottomrule
\end{tabularx}%
}


\begin{center}\rule{0.5\linewidth}{0.5pt}\end{center}


\subsecao{6.5 Acompanhamento da execução}

Para assegurar que a melhoria incorporada ao manual continue sendo praticada de forma sustentável nas próximas iterações e ao longo das Macroentregas 2 e 3, a CineSoft Solutions estabeleceu a seguinte sistemática de verificação operacional:

\begin{enumerate}
\def\labelenumi{\arabic{enumi}.}
\tightlist
\item
  \textbf{Quem realiza a conferência:} \textbf{Kayque Vicente Silva} (Gerente de Projeto \& QA Lead), em conjunto com \textbf{Gabriel Henrique Lourenço Dias} (Tech Lead).
\item
  \textbf{Quando ocorre a verificação:} Semanalmente, durante a cerimônia de \emph{Sprint Planning / Review}, e continuamente a cada abertura de Pull Request no repositório do GitHub.
\item
  \textbf{Qual registro e medida são observados:}

  \begin{itemize}
  \tightlist
  \item
    \textbf{Indicador de Conformidade do Quadro Kanban:} Percentual de cartões ativos nas colunas \emph{Sprint Backlog} e \emph{In Progress} do GitHub Projects que possuem simultaneamente \textbf{responsável atribuído (\texttt{Assignee})}, \textbf{prazo estipulado (\texttt{Iteration\ /\ Milestone})} e \textbf{requisitos vinculados} (Meta: \textbf{100\% dos cartões conformes}).
  \item
    \textbf{Registro de Auditoria:} Caso seja identificado qualquer cartão sem responsável ou acima do limite de 48h sem atualização, a pendência é registrada e corrigida imediatamente na abertura da reunião semanal, acionando o protocolo de pareamento descrito na seção 4.1.
  \end{itemize}
\end{enumerate}

\begin{center}\rule{0.5\linewidth}{0.5pt}\end{center}


\subsubsecao{Referências Normativas do Capítulo}

\begin{itemize}
\tightlist
\item
  \textbf{SOFTEX.} \emph{Guia Geral MPS de Software: 2024 (MR-MPS-SW)}. Associação para Promoção da Excelência do Software Brasileiro. Disponível em: https://softex.br/download/guia-geral-mps-de-software2024/.
\item
  \textbf{ABNT.} \emph{NBR ISO/IEC/IEEE 12207:2021 --- Engenharia de sistemas e software --- Processos de ciclo de vida de software}. Rio de Janeiro: ABNT, 2021.
\end{itemize}
